\documentclass[11pt,a4paper]{article}

\usepackage[utf8]{inputenc}
\usepackage[T1]{fontenc}
\usepackage{amsmath,amssymb,amsthm}
\usepackage{graphicx}
\usepackage{booktabs}
\usepackage{hyperref}
\usepackage{xcolor}
\usepackage[margin=25mm]{geometry}
\usepackage{xeCJK}
\setCJKmainfont{Hiragino Mincho ProN}
\setCJKsansfont{Hiragino Sans}

\newtheorem{theorem}{定理}
\newtheorem{proposition}{命題}
\newtheorem{observation}{観察}
\newtheorem{conjecture}{予想}
\newtheorem{identity}{恒等式}
\newtheorem{corollary}{系}

\title{Eisenstein Fourier 係数としての物理定数:\\
観測的恒等式と Wheeler--DeWitt による $\Omega_\Lambda = 2\pi/9$ の導出}

\author{Wright Brothers Research Program}
\date{2026年4月}

\begin{document}

\maketitle

\begin{abstract}
本稿は観測される無次元物理定数と正規化 Eisenstein 級数の Fourier
係数を結ぶ算術的恒等式を報告する．具体的には，微細構造定数の
主要項 $\alpha^{-1} = 2/|B_2| + 4/|B_4| + \ldots = 132 + \ldots$ が
$(E_4 - E_2)/2$ の $q^1$ 係数に等しい．Wright Brothers の ``Bernoulli
ladder'' $L_k = 2k/|B_{2k}| \in \{12, 120, 252, 240, 132\}$ は
正規化 Eisenstein 級数の $q^1$ 係数の絶対値の半分に厳密に等しい．
ダークエネルギー密度パラメータは関数等式の双対形式
$\Omega_\Lambda = 16|\zeta(-1)|\zeta(2)/\pi = 2\pi/9 \approx 0.6981$
に書き換えられる．この $\Omega_\Lambda$ 値は量子宇宙論からも独立に
得られる: 最小超空間 Wheeler--DeWitt 理論において，宇宙波動関数
$\Psi(a)$ は時間的 Dirichlet--Neumann 境界条件 (Big Bang で Dirichlet,
de Sitter 漸近で Neumann) を満たし，その 1 次元零点エネルギーは
$\zeta$-正則化で $\zeta_{\neg 2}(-1) = +1/12$ に評価される．
Cohen--Kaplan--Nelson ホログラフィック境界と組み合わせると，
Planck 2018 観測値 $0.6847 \pm 0.0073$ と 2\% ($1.8\sigma$) の精度で
一致する $\Omega_\Lambda$ が得られる．本稿ではこれらの恒等式を
既存の保型形式物理 (ヘテロティック弦，moonshine, 位相弦) および
Kreimer--Connes の摂動 QED 繰り込み Hopf 代数と関連付けて論じる．
観察された事実は，特定の無次元物理定数が整数引数における $\zeta$
特殊値を含む保型形式構成の評価として表現可能であること---すなわち
Deligne--Beilinson 周期予想の特定の物理的実現---を示唆する．
限界 (ほとんどの恒等式の厳密第一原理導出の不在を含む) を明示的に
列挙し，将来の宇宙論サーベイ (Euclid, DESI, LSST) と精密粒子測定
(Belle II, ミューオン $g-2$) による具体的検証経路を同定する．
\end{abstract}

\section{序}

理論物理学の長年の二大問題は，無次元基本定数 (例えば微細構造定数
$\alpha^{-1} = 137.035999\ldots$)~\cite{Mohr2016} の起源と，宇宙定数の
数値 ($\rho_\Lambda \approx 5.24 \times 10^{-10}$ J/m$^3$，つまり
$\Omega_\Lambda = 0.6847 \pm 0.0073$ (Planck 2018)~\cite{Planck2018}) で
ある．前者は現在の理論から予測的導出が得られておらず，後者は素朴な
QFT 推定値より $10^{123}$ 桁小さい~\cite{Weinberg1989}．

本研究はこれらの定数 (および異常ミューオン磁気モーメント
$\Delta a_\mu \approx 251 \times 10^{-11}$~\cite{MuonG2FNAL} を含む他の定数)
と正規化 Eisenstein 級数 $E_{2k}(\tau)$ およびその関連保型形式
対象の Fourier 係数とを結びつける観察的恒等式を報告する．保型形式は
数理物理学に長く現れてきた---ヘテロティック弦理論~\cite{GSW1987},
モンスター moonshine~\cite{Borcherds1992}, Mathieu
moonshine~\cite{EOT2010}, 位相弦不変量, ブラックホール状態
数え上げなど---が，これらの系統的な\emph{低エネルギー無次元定数の
数値の決定への利用}は文献上ほぼ存在しない．本稿はそのような
プログラムを提案する．

第二の主結果は $\Omega_\Lambda$ の量子宇宙論的議論である．最小超空間
Wheeler--DeWitt 波動関数 $\Psi(a)$ が Big Bang ($a = 0$) での Dirichlet
境界条件と未来無限遠 ($a \to \infty$) での Neumann 条件を満たせば，
対応する 1 次元零点エネルギーは $\zeta_{\neg 2}(-1) = +1/12$ で
正則化される．Cohen--Kaplan--Nelson ホログラフィック
境界~\cite{CKN1999}と組み合わせると，
\begin{equation}
\Omega_\Lambda = \frac{2\pi}{9} \approx 0.6981
\end{equation}
がパラメータ自由な予測として得られ，Planck 2018 と $1.8\sigma$ の
水準で整合する．

本手法は宇宙定数問題への正準的 ``弦ランドスケープ'' 対処
法~\cite{Susskind2003}と異なり，多宇宙上の人間原理的選択でなく
特定の数値を産出する．

\subsection{論文の構成}

第 \ref{sec:notation}節で記号を定める．第 \ref{sec:main}節で主要な
恒等式を述べる．第 \ref{sec:wdw}節で $\Omega_\Lambda = 2\pi/9$ の
Wheeler--DeWitt 時間的 DN 導出を提示する．第 \ref{sec:context}節で
既存の保型形式物理における文脈を配置する．第 \ref{sec:predictions}節で
具体的検証可能予測を列挙する．第 \ref{sec:limitations}節で限界を
honest に評価する．

\section{記号と準備}
\label{sec:notation}

標準的に正規化された Eisenstein 級数 (重さ $2k \geq 4$):
\begin{equation}
E_{2k}(\tau) = 1 - \frac{4k}{B_{2k}}\sum_{n=1}^\infty \sigma_{2k-1}(n)\,q^n,
\quad q = e^{2\pi i\tau}
\end{equation}
ここで $B_{2k}$ はベルヌーイ数, $\sigma_s(n) = \sum_{d|n}d^s$．
$E_{2k}$ の $q^1$ Fourier 係数は $-4k/B_{2k}$．

準保型 Eisenstein 級数 $E_2(\tau)$ は同じ式を $k=1$ で用いる
(すなわち $q^1$ 係数は $-24$)．これは以下により保型不変性を
破る:
\begin{equation}
E_2\bigl(\tfrac{a\tau+b}{c\tau+d}\bigr) = (c\tau+d)^2 E_2(\tau) - \frac{6ic(c\tau+d)}{\pi},
\end{equation}
保型完成は $E_2^*(\tau) = E_2(\tau) - 3/(\pi\,\mathrm{Im}\,\tau)$．

Riemann ゼータ関数は
\begin{equation}
\zeta(1-2k) = -\frac{B_{2k}}{2k}, \qquad \zeta(2k) = \frac{(-1)^{k+1}(2\pi)^{2k}B_{2k}}{2(2k)!},
\end{equation}
を満たし，したがって $|\zeta(-1)| = 1/12$, $|\zeta(-3)| = 1/120$,
$|\zeta(-5)| = 1/252$ 等．

\emph{Wright Brothers Bernoulli ladder} を次のように定義する:
\begin{equation}
L_k \equiv \frac{2k}{|B_{2k}|} = \frac{1}{|\zeta(1-2k)|}.
\end{equation}
最初のエントリーは
\begin{equation}
(L_1, L_2, L_3, L_4, L_5) = (12,\,120,\,252,\,240,\,132).
\end{equation}
$k=6$ で，irregular prime $691$ が $B_{12}$ の numerator を割る
~\cite{Kummer1850}ため $L_6 = 32760/691$ は非整数．

\section{主要な観察}
\label{sec:main}

\subsection{ladder--Eisenstein 恒等式}

\begin{identity}[ladder--Eisenstein]
\label{id:ladder}
$k \geq 1$ に対して，
\begin{equation}
L_k = \frac{2k}{|B_{2k}|} = \frac{1}{2}\bigl|E_{2k}\text{ の }q^1\text{ 係数}\bigr|.
\end{equation}
\end{identity}

\begin{proof}
$E_{2k}$ の $q^1$ 係数は $-4k/B_{2k}$．絶対値を取り 2 で割ると
$2k/|B_{2k}| = L_k$ が得られる．
\end{proof}

数学的には初等的だが，この恒等式は Wright Brothers 枠組み全体で
登場する Bernoulli ladder エントリーに対する系統的な保型形式源を
与える．

\subsection{$\alpha^{-1}$ 主要項と $(E_4 - E_2)/2$}

\begin{theorem}[$\alpha^{-1}$ 保型恒等式]
\label{th:alpha}
\begin{equation}
\frac{E_4(\tau) - E_2(\tau)}{2} = 1 + 132\, q + O(q^2)
\end{equation}
ここで $132 = 2/|B_2| + 4/|B_4|$ は Wright Brothers の $\alpha^{-1}$
公式の主要算術項に等しい．
\end{theorem}

\begin{proof}
$E_4 = 1 + 240q + O(q^2)$ かつ $E_2 = 1 - 24q + O(q^2)$ であるから，
$(E_4 - E_2)/2 = 1 + (240 - (-24))/2 \cdot q + O(q^2) = 1 + 132q + O(q^2)$．
等式 $132 = 12 + 120 = 2/|B_2| + 4/|B_4|$ は $|B_2| = 1/6$ かつ
$|B_4| = 1/30$ から直ちに従う．
\end{proof}

\subsection{ミューオン $g-2$ 異常: 三重表現}

実験的に観測されたミューオン異常磁気モーメント差
$\Delta a_\mu = (251 \pm 59)\times 10^{-11}$~\cite{MuonG2FNAL, BMW2021}は，
数値 $252$ に対する三つの独立した「保型形式的」表現を許す:

\begin{observation}[$252$ の三重表現]
\label{obs:muon}
\begin{enumerate}
\item (Bernoulli ladder): $L_3 = 6/|B_6| = 252$
\item (Eisenstein 差):
$(E_8 - E_2)/2 = 1 + 252\,q + O(q^2)$
\item (Ramanujan 判別式): $\tau(3) = 252$, ここで
$\Delta(\tau) = \sum \tau(n)q^n$ は重さ 12 の cusp form
\end{enumerate}
\end{observation}

\begin{proof}
(1) $L_3 = 6/|1/42| = 252$．
(2) $E_8 = 1 + 480q + O(q^2)$ (重さ 8 で $E_8 = E_4^2$) なので，
$(E_8 - E_2)/2 = 1 + (480+24)/2 \cdot q + O(q^2) = 1 + 252q + O(q^2)$．
(3) $\tau(3) = 252$ は Ramanujan tau 関数の古典値．
\end{proof}

数値的一致 $\Delta a_\mu \approx 252$ (単位 $10^{-11}$) は先行研究で
指摘されている~\cite{wb-prior}．新しい観察は，この数が同時に単一
Eisenstein ladder エントリー，Eisenstein 差，cusp form 係数の
三つの独立した保型形式対象に等しいことである．

\subsection{ladder の非自明な加法恒等式}

\begin{identity}[ladder 加法関係]
\label{id:additive}
Wright Brothers ladder 内で，以下の非自明な恒等式が成立する:
\begin{align}
L_1 + L_2 &= L_5 &\text{(すなわち $12 + 120 = 132$)} \\
L_1 + L_4 &= L_3 &\text{(すなわち $12 + 240 = 252$)} \\
L_2 + L_5 &= L_3 &\text{(すなわち $120 + 132 = 252$)}
\end{align}
\end{identity}

これらは $|B_2|^{-1} + 2|B_4|^{-1} = 5|B_{10}|^{-1}$ (つまり
$6 + 60 = 66$) のような非自明な Bernoulli 関係を表現する．
これは一旦述べられれば初等的だが，標準的な Bernoulli 漸化式からの
明白な導出を持たない．特定の和のもとでの ladder の閉性は，さらに
調査すべき隠れた構造を示唆する．

\subsection{関数等式双対としての $\Omega_\Lambda$}

\begin{identity}[$\Omega_\Lambda$ の関数等式形式]
\label{id:omega}
\begin{equation}
\Omega_\Lambda = \frac{2\pi}{9} = 8\pi B_2^2 = -\frac{16\,\zeta(-1)\,\zeta(2)}{\pi}
= \frac{16\,|\zeta(-1)|\,\zeta(2)}{\pi}.
\end{equation}
\end{identity}

\begin{proof}
$B_2 = 1/6$ なので $8\pi B_2^2 = 8\pi/36 = 2\pi/9$．ゼータ形式に対して:
$\zeta(-1) = -1/12$, $\zeta(2) = \pi^2/6$ であるから
$-16\zeta(-1)\zeta(2)/\pi = 16 \cdot (1/12)(\pi^2/6)/\pi = 16\pi/72 = 2\pi/9$．
\end{proof}

この形式は，負の整数での $\zeta$ (Bernoulli 経由で algebraic な
$\zeta(-1)$) と正の整数での $\zeta$ ($\pi^2$ を含む周期
$\zeta(2) = \pi^2/6$) を，明示的な $1/\pi$ 正則化因子と共に結合する
点で注目される．このような組合せは Riemann 関数等式の両辺を用いる．

\subsection{$E_2$ 準保型性との関連}

恒等式 \ref{id:omega}の $1/\pi$ 因子は $E_2$ の保型完成と intriguingly
平行している:
\begin{equation}
E_2^*(\tau) = E_2(\tau) - \frac{3}{\pi\,\mathrm{Im}\,\tau},
\end{equation}
ここで $3/(\pi y)$ が準保型補正である．特定の点での $E_2^*$ の評価から
$\Omega_\Lambda$ を直接導くことはまだできないが，共通の $1/\pi$ 構造は
接続を示唆する．既知の特殊値は $E_2(i) = 3/\pi$~\cite{Serre1973}で，
これから $E_2^*(i) = 0$．$\Omega_\Lambda = 2\pi/9 \neq (E_2(i))^{-1}$ なので，
自明な同一化は存在しない．

\section{Wheeler--DeWitt による $\Omega_\Lambda$ の導出}
\label{sec:wdw}

ここで $\Omega_\Lambda = 2\pi/9$ を量子宇宙論から独立に導出する．

\subsection{最小超空間 Wheeler--DeWitt 方程式}

Wheeler--DeWitt 方程式~\cite{DeWitt1967,Wheeler1968}は，インフラトン
$\phi$ を含む平坦 FRW 宇宙の最小超空間で，波動関数 $\Psi(a,\phi)$
に対する 1 次元 Schr\"odinger 様の方程式に簡約される:
\begin{equation}
\left[-\frac{\hbar^2}{2}\frac{\partial^2}{\partial a^2} + V(a) + \hat H_\phi\right]\Psi(a,\phi) = 0.
\end{equation}
ポテンシャル $V(a)$ は曲率，宇宙定数，物質寄与を統合
する~\cite{Halliwell1990}．

\subsection{時間的 Dirichlet--Neumann 境界条件}

$\Psi(a)$ が以下の境界条件を満たすと提案する:

\begin{proposition}[時間的 DN 境界]
\begin{align}
\Psi(a=0) &= 0 \quad\text{(Big Bang での Dirichlet)},\\
\partial_a\Psi|_{a\to\infty} &\text{ 自由} \quad\text{(de Sitter 漸近での Neumann)}.
\end{align}
\end{proposition}

正当化: (i) $a = 0$ は Big Bang 特異点に対応し，ここで古典幾何が破綻
する．要請 $\Psi(0) = 0$ は Hartle--Hawking の無境界提案
~\cite{HartleHawking1983}の Lorentzian 符号解釈および Vilenkin の
トンネル提案~\cite{Vilenkin1982}と整合する．(ii) 遠未来では宇宙は
定数 $\Lambda$ の de Sitter に漸近し，WKB 近似で $\Psi(a) \sim
a^{-3/2}e^{\pm i a^3\sqrt\Lambda/(3\hbar)}$ となり，Neumann 様 (自由振動)
境界条件と整合する．

\subsection{1D DN 零点エネルギー}

長さ $L$ で Dirichlet--Neumann 境界条件をもつ 1 次元調和空洞は
モード $k_n = (2n-1)\pi/(2L)$, $n = 1, 2, \ldots$ を持つ．ゼータ関数で
正則化された零点エネルギーは
\begin{equation}
E_0^{\rm DN} = \frac{\hbar c\pi}{4L}\sum_{n=1}^\infty(2n-1)_{\text{reg}}
= \frac{\hbar c\pi}{4L}\cdot\zeta_{\neg 2}(-1)
\end{equation}
ここで
\begin{equation}
\zeta_{\neg 2}(s) \equiv \sum_{\text{奇数 }n \geq 1}n^{-s} = (1 - 2^{-s})\zeta(s).
\end{equation}
$s = -1$ で:
\begin{equation}
\zeta_{\neg 2}(-1) = (1-2)\zeta(-1) = (-1)(-1/12) = +\frac{1}{12}.
\end{equation}

正の符号が本質的である．対照的に NN または DD 空洞は $\zeta(-1) = -1/12$
を用い，引力的 (負) Casimir エネルギーを与える．DN の場合は正の
Casimir エネルギーを生じ，ダークエネルギー ($\rho_\Lambda > 0$) に
適している．

\subsection{Cohen--Kaplan--Nelson 境界}

CKN 境界~\cite{CKN1999}は，サイズ $L$ の領域の真空エネルギーが
\begin{equation}
\rho_{\rm vac}(L) \lesssim M_P^2/L^2 = \rho_P\bigl(\ell_P/L\bigr)^2.
\end{equation}
を満たすと述べる．$L = \ell_H$ (Hubble 長) の観測可能宇宙に適用し,
DN 構造から係数 $|\zeta_{\neg 2}(-1)| = 1/12$ での飽和を仮定すると，
\begin{equation}
\rho_\Lambda = \frac{1}{12}\cdot\rho_P\cdot\bigl(\ell_P/\ell_H\bigr)^2.
\end{equation}

\subsection{$\Omega_\Lambda = 2\pi/9$ の導出}

$\rho_P = \hbar c/\ell_P^4$ と $\ell_P^2 = \hbar G/c^3$ を用いて:
\begin{equation}
\rho_\Lambda = \frac{1}{12}\cdot\frac{c^4}{G\ell_H^2}.
\end{equation}
観測的定義は
\begin{equation}
\rho_\Lambda = \Omega_\Lambda\cdot\rho_c = \Omega_\Lambda\cdot\frac{3c^4}{8\pi G\ell_H^2},
\end{equation}
($\ell_H = c/H_0$ を用いる) を与える．両式を等置:
\begin{equation}
\frac{1}{12} = \frac{3\Omega_\Lambda}{8\pi}
\quad\Longrightarrow\quad
\boxed{\Omega_\Lambda = \frac{2\pi}{9}}.
\end{equation}

数値的には $2\pi/9 = 0.69813$ で，Planck 2018 観測値
$0.6847 \pm 0.0073$ と $1.8\sigma$ (2.0\%) で整合する．特筆すべきは，
予測が $H_0$ に依存しないことである．$\ell_H$ が無次元比で相殺する
ためである．これは Planck と SH0ES の $H_0$ 値の緊張が $\Omega_\Lambda$
予測に影響しないことを意味する．

\section{文脈: 保型形式物理}
\label{sec:context}

保型形式は理論物理学で長く重要であった．最も顕著な応用:

\begin{itemize}
\item \textbf{ヘテロティック弦理論}: $E_8 \times E_8$ および $SO(32)$
ヘテロティック弦の分配関数は偶ユニモジュラー格子の Eisenstein 級数と
テータ関数を含む~\cite{GSW1987}．特に $E_8$ 格子のテータ関数は
$E_4$ に等しい: $\theta_{E_8}(\tau) = E_4(\tau)$~\cite{Serre1973}．

\item \textbf{Monstrous moonshine}: $j$ 関数の Fourier 係数は
Monster 群表現の次元であり，Borcherds により頂点作用素代数を
介して証明された~\cite{Borcherds1992}．

\item \textbf{Mathieu moonshine}: K3 面の楕円種数は mock modular
form を含み，その Fourier 係数は $M_{24}$ 表現の次元に分解される
~\cite{EOT2010,Gannon2016}．

\item \textbf{位相弦}: Calabi--Yau 3-fold 上の分配関数は
Gopakumar--Vafa 不変量を encode する準保型形式を生成する
~\cite{GV1998,BCOV1994}．

\item \textbf{ブラックホール microstate}: $\mathcal{N}=4$ dyon 退化度は
Igusa cusp form~\cite{DMVV1997}および mock Jacobi form~\cite{Dabholkar2012}
の Fourier 係数である．

\item \textbf{$\mathcal{N}=4$ SYM S 双対}: 分配関数は
$SL_2(\mathbb{Z})$ 不変である~\cite{VafaWitten1994}．

\item \textbf{Kreimer--Connes 繰り込み}: 摂動 QED Feynman 積分は
multiple zeta 値および polylogarithm に評価される
~\cite{Kreimer1997,ConnesKreimer2000,Brown2015}．
\end{itemize}

本研究はこれらに対して相補的である．既存の応用は主に高次元弦理論,
トポロジカル部門, または形式的な BPS 状態数え上げに関するもの．
一方，第 \ref{sec:main}節の恒等式は代わりに，4 次元物理世界の
低エネルギー無次元定数の\emph{数値}を保型形式対象に結び付ける．

これは Kreimer--Connes 枠組みに最も近い．そこでは摂動 QED 計算が
arithmetic 定数を生成する．しかし Kreimer--Connes は
$\alpha^{-1} = 137.036$ や $\Omega_\Lambda = 0.698$ のような特定の値を
予測しない．ループ単位で計算する枠組みを提供するが，値自体は
通常の場の量子論から emerge する．本研究は代わりに，ループ展開を
迂回して閉形式の arithmetic 表現を与える直接的な保型形式的
``辞書'' を提案する．

\subsection{Deligne--Beilinson 周期予想との関連}

Deligne 予想~\cite{Deligne1979}および Beilinson 予想~\cite{Beilinson1984}
は，モチーフの臨界 $L$ 値が周期と代数的数の積として factor 化すると
予測する．形式
\begin{equation}
\Omega_\Lambda = \frac{16|\zeta(-1)|\zeta(2)}{\pi}
\end{equation}
はそのような factor 化を想起させる: $\zeta(-1)$ は有理数 (algebraic),
$\zeta(2) = \pi^2/6$ は周期であり，$1/\pi$ は次元因子である．
Wheeler--DeWitt 導出が正しければ，$\Omega_\Lambda$ は未同定モチーフに
対する Beilinson 予想の物理的実例となる．

\section{予測と検証}
\label{sec:predictions}

本枠組みの検証可能予測を列挙する．

\subsection{宇宙論サーベイ}

\begin{itemize}
\item \textbf{$\Omega_\Lambda = 2\pi/9$}: 将来の宇宙論サーベイ
(Euclid~\cite{Euclid}, DESI~\cite{DESI}, LSST~\cite{LSST}) は
$\Omega_\Lambda$ を $\sim 0.5\%$ の水準で測定する．
$2\pi/9 = 0.69813$ 対現在の best-fit $0.6847$ の決定的検証．

\item \textbf{状態方程式 $w = -1$ exact}: 導出は $\rho_\Lambda$ が定数と
仮定し，$w_{\rm DE} = -1$ を時間変化なしで予測する．最近の DESI 2024
による $w(z) \neq -1$ の hint~\cite{DESI2024}が確認されれば本予測は
棄却される．
\end{itemize}

\subsection{粒子物理}

\begin{itemize}
\item \textbf{ミューオン $g-2$ 精密}: 同定
$\Delta a_\mu \leftrightarrow \tau(3) = 252$ は精密な規格化を示唆する．
現在の Fermilab+BNL 合算値 $(251 \pm 59)\times 10^{-11}$~\cite{MuonG2FNAL}
は不確かさ内．

\item \textbf{タウ $g-2$}: 世代-$\tau(n)$ 対応が成立すれば，
$\Delta a_\tau$ は $\tau(5) = 4830$ に比例するはず．現在 $10^{-2}$
精度では検証不能．Belle II と将来コライダーで $10^{-4}$ 精度に
達するかもしれない．
\end{itemize}

\subsection{実験室凝縮系}

\begin{itemize}
\item \textbf{ダークエネルギー長スケール}: 予測
$\rho_\Lambda \approx 5.3 \times 10^{-10}$ J/m$^3$ は特徴長
$\ell_\Lambda = (\hbar c/\rho_\Lambda)^{1/4} \approx 88\,\mu$m に対応する．
このスケールでの Casimir 力および関連する真空効果の実験室試験
(水晶バルク音響共振器の厚さ系列) は相転移の signature を明らかに
する可能性がある~\cite{wb-bawguide}．

\item \textbf{Dirichlet--Neumann 境界 Casimir}: 枠組みは混合 DN 構造
を持つ境界での Casimir 物理の変化を予測する．商用 solidly-mounted
resonator (SMR) はそのような境界を提供し，この文脈で分析されて
いない．
\end{itemize}

\section{限界}
\label{sec:limitations}

本研究の限界を honest に列挙する:

\begin{enumerate}
\item \textbf{恒等式の数学的自明性}: 恒等式 \ref{id:ladder}および
定理 \ref{th:alpha}のような恒等式はベルヌーイ数および Eisenstein
級数の定義から直ちに従う．新規性は数学的内容自体ではなく，
\emph{物理への接続}にある．

\item \textbf{ほとんどの観察の post-hoc 性質}: 接続
$\alpha^{-1} \leftrightarrow (E_4 - E_2)/2$ と
$\Delta a_\mu \leftrightarrow \tau(3)$ は ab initio に予測される
のでなく事後的に観察される．Look-elsewhere effect は完全には制御
できない．

\item \textbf{CKN 係数}: $\Omega_\Lambda$ の Wheeler--DeWitt 導出は
係数 $1/12$ での CKN 境界の飽和を仮定する．DN モード正則化を
介して正当化するが，この係数の厳密第一原理導出はまだ得られていない．

\item \textbf{厳密 Dirichlet 問題}: モード分解での $a = 0$ での厳密
Dirichlet 境界は非正規化可能な Bogoliubov 係数を導く．
smoothing parameter を導入する必要がある~\cite{wb-low-ell}．この
smoothing の正確な形は一意でない．

\item \textbf{quadrupole 異常}: DN 境界からの低 $\ell$ CMB 抑制予測は
観測された $\ell = 2$ quadrupole 抑制 ($\approx 40\%$) を定量的に
1 桁過小評価する．単純な smoothed Dirichlet は多くても $\sim 10\%$
抑制しか生成しない~\cite{wb-low-ell}．

\item \textbf{質量階層と混合}: 枠組みはフェルミオン質量，Yukawa 結合,
CKM/PMNS 混合を扱わない．これらは $10^{-7}$ 以上の観測精度を持つが，
Bernoulli ladder 型の arithmetic 表現の類似物を持たない．

\item \textbf{ラグランジアンからの厳密導出}: 既知の QFT または重力
ラグランジアンから観察された恒等式を第一原理から導出できない．
Kreimer--Connes 枠組みがそのような導出を可能にするツールを提供する
可能性があるが，これは将来の仕事のままである．
\end{enumerate}

\section{議論}

ここで報告される観察は，特定の無次元物理定数---特に $\alpha^{-1}$,
$\Omega_\Lambda$, $\Delta a_\mu$---が Riemann $\zeta$ 特殊値の保型
形式評価または関数等式の組合せとして表現可能であることを示唆する．
これらの表現のステータスは初等恒等式 (ladder--Eisenstein 対応)
から speculative な予想 (世代--Ramanujan-$\tau$ 仮説) まで多様で
ある．

$\Omega_\Lambda = 2\pi/9$ の Wheeler--DeWitt 導出は，量子宇宙論と
保型形式 arithmetic の間でこのプログラムにおける最も具体的な link
を提供する．将来の宇宙論サーベイにより $0.5\%$ の水準で反証可能である．

もしこれらの接続が偶然でなく構造的であれば，数論幾何，特に
Langlands プログラムおよび Beilinson 予想の対象物が，物理定数の
決定において深い役割を果たすことを示唆する．``Wright Brothers
alphabet''---ベルヌーイ数，$\zeta$ 特殊値，Ramanujan $\tau$, および
関連する保型形式データ---は無次元パラメータに対する系統的言語を
構成するであろう．

もし偶然であれば，プログラムは特定の予測 (特に $\Omega_\Lambda
= 2\pi/9$ の $< 0.5\%$ 精度およびタウ $g-2$) で失敗し，観察は
数値的偶然として片付けられる．どちらの結果も科学的に生産的である．

\section{結論}

本研究は観測される無次元物理定数を正規化 Eisenstein 級数の Fourier
係数および Riemann $\zeta$ 特殊値に結ぶ算術的恒等式をいくつか報告
した．最も具体的な結果は:
\begin{enumerate}
\item[(i)] Wright Brothers Bernoulli ladder $L_k = 2k/|B_{2k}|$ は
$E_{2k}$ の Eisenstein $q^1$ 係数の絶対値の半分に厳密に等しい．
\item[(ii)] Wright Brothers の $\alpha^{-1}$ 公式の主要算術部分
(132 に等しい) は $(E_4 - E_2)/2$ の $q^1$ 係数と一致する．
\item[(iii)] ミューオン異常磁気モーメント
$\Delta a_\mu \approx 252 \times 10^{-11}$ は三重保型形式表現を持つ:
$L_3$ として，$(E_8 - E_2)/2|_{q^1}$ として，Ramanujan 判別式係数
$\tau(3)$ として．
\item[(iv)] ダークエネルギー密度パラメータ
$\Omega_\Lambda = 2\pi/9 \approx 0.6981$ は関数等式双対形式
$\Omega_\Lambda = 16|\zeta(-1)|\zeta(2)/\pi$ を許す．
\item[(v)] 同じ $\Omega_\Lambda = 2\pi/9$ は，時間的 Dirichlet--Neumann
境界条件を持つ Wheeler--DeWitt 最小超空間議論と Cohen--Kaplan--Nelson
ホログラフィック境界の組合せから独立に生じる．
\end{enumerate}

これらの結果は将来の宇宙論サーベイで $0.5\%$ の水準で検証可能である．
成り立てば，ダークエネルギーと量子宇宙論の文脈内で Deligne--Beilinson
周期予想の特定の物理的実現を提供する．

\section*{謝辞}

本研究は Wright Brothers Research Program，独立した研究プロジェクトの
一部である．著者は Feynman 振幅の arithmetic 理解への D. Kreimer,
A. Connes, M. Marcolli, F. Brown の基礎的貢献，および保型形式に
関する R. Borcherds, J.-P. Serre 他の仕事を謝意をもって認める．

\begin{thebibliography}{99}

\bibitem{Mohr2016}
P.~J.~Mohr, D.~B.~Newell, B.~N.~Taylor, Rev.\ Mod.\ Phys.\ \textbf{88}, 035009 (2016).

\bibitem{Planck2018}
Planck Collaboration, Astron.\ Astrophys.\ \textbf{641}, A6 (2020).

\bibitem{Weinberg1989}
S.~Weinberg, Rev.\ Mod.\ Phys.\ \textbf{61}, 1 (1989).

\bibitem{MuonG2FNAL}
D.~P.~Aguillard \textit{et al.}\ (Muon $g-2$ Collaboration),
Phys.\ Rev.\ Lett.\ \textbf{131}, 161802 (2023).

\bibitem{BMW2021}
S.~Borsanyi \textit{et al.}\ (BMW Collaboration),
Nature \textbf{593}, 51 (2021).

\bibitem{CKN1999}
A.~G.~Cohen, D.~B.~Kaplan, A.~E.~Nelson,
Phys.\ Rev.\ Lett.\ \textbf{82}, 4971 (1999).

\bibitem{DeWitt1967}
B.~S.~DeWitt, Phys.\ Rev.\ \textbf{160}, 1113 (1967).

\bibitem{Wheeler1968}
J.~A.~Wheeler, in \textit{Battelle Rencontres},
eds.\ C.~M.~DeWitt and J.~A.~Wheeler (Benjamin, 1968).

\bibitem{Halliwell1990}
J.~J.~Halliwell, in \textit{Quantum Cosmology and Baby Universes},
eds.\ S.~Coleman \textit{et al.}\ (World Scientific, 1991).

\bibitem{HartleHawking1983}
J.~B.~Hartle, S.~W.~Hawking, Phys.\ Rev.\ D \textbf{28}, 2960 (1983).

\bibitem{Vilenkin1982}
A.~Vilenkin, Phys.\ Lett.\ B \textbf{117}, 25 (1982).

\bibitem{Susskind2003}
L.~Susskind, hep-th/0302219.

\bibitem{GSW1987}
M.~B.~Green, J.~H.~Schwarz, E.~Witten, \textit{Superstring Theory}
(Cambridge Univ.\ Press, 1987).

\bibitem{Borcherds1992}
R.~Borcherds, Invent.\ Math.\ \textbf{109}, 405 (1992).

\bibitem{EOT2010}
T.~Eguchi, H.~Ooguri, Y.~Tachikawa, Exp.\ Math.\ \textbf{20}, 91 (2011).

\bibitem{Gannon2016}
T.~Gannon, Adv.\ Math.\ \textbf{301}, 322 (2016).

\bibitem{GV1998}
R.~Gopakumar, C.~Vafa, hep-th/9809187 and hep-th/9812127 (1998).

\bibitem{BCOV1994}
M.~Bershadsky, S.~Cecotti, H.~Ooguri, C.~Vafa,
Commun.\ Math.\ Phys.\ \textbf{165}, 311 (1994).

\bibitem{DMVV1997}
R.~Dijkgraaf, G.~Moore, E.~Verlinde, H.~Verlinde,
Commun.\ Math.\ Phys.\ \textbf{185}, 197 (1997).

\bibitem{Dabholkar2012}
A.~Dabholkar, S.~Murthy, D.~Zagier, arXiv:1208.4074.

\bibitem{VafaWitten1994}
C.~Vafa, E.~Witten, Nucl.\ Phys.\ B \textbf{431}, 3 (1994).

\bibitem{Kreimer1997}
D.~Kreimer, Adv.\ Theor.\ Math.\ Phys.\ \textbf{2}, 303 (1998),
q-alg/9707029.

\bibitem{ConnesKreimer2000}
A.~Connes, D.~Kreimer, Commun.\ Math.\ Phys.\ \textbf{210}, 249 (2000).

\bibitem{Brown2015}
F.~Brown, Ann.\ Math.\ \textbf{181}, 949 (2015).

\bibitem{Serre1973}
J.-P.~Serre, \textit{A Course in Arithmetic} (Springer, 1973).

\bibitem{Kummer1850}
E.~E.~Kummer, J.\ Reine Angew.\ Math.\ \textbf{40}, 131 (1850).

\bibitem{Deligne1979}
P.~Deligne, Proc.\ Symp.\ Pure Math.\ \textbf{33.2}, 313 (1979).

\bibitem{Beilinson1984}
A.~Beilinson, J.\ Soviet Math.\ \textbf{30}, 2036 (1985).

\bibitem{Euclid}
R.~Laureijs \textit{et al.}, arXiv:1110.3193.

\bibitem{DESI}
DESI Collaboration, arXiv:1611.00036.

\bibitem{LSST}
LSST Science Collaborations, arXiv:0912.0201.

\bibitem{DESI2024}
DESI Collaboration, arXiv:2404.03002 (2024).

\bibitem{wb-prior}
Wright Brothers Research Program, \texttt{b2\_unification\_deeper\_ja.tex} (2026).

\bibitem{wb-bawguide}
Wright Brothers Research Program, \texttt{guide\_baw\_quartz\_beginners\_ja.tex} (2026).

\bibitem{wb-low-ell}
Wright Brothers Research Program, \texttt{low\_ell\_cmb\_dn\_ja.tex} (2026).

\end{thebibliography}

\end{document}
